\documentclass[fleqn]{article}
\usepackage{amsmath}
\usepackage{amssymb}
\usepackage{array}
\setlength{\mathindent}{1.5em}

\begin{document}

\begin{center}
{\LARGE \textbf{Radix-4 Booth}}\\[4pt]
{\large signed digits from an overlapping 3-bit window of $B$}
\end{center}

\vspace{3em}

\noindent\textbf{Step 1 --- the number}

\noindent $B$ is an $n$-bit two's complement value ($n$ even). The top bit counts negative:
\[
B \;=\; -\,2^{n-1} B^{\,n-1} \;+\; \sum_{j=0}^{n-2} 2^{j} \, B^{\,j}
\]

\vspace{0.5em}

\noindent\textbf{Step 2 --- the digit}

\noindent Each digit reads a 3-bit window; neighbouring digits share one bit ($B^{\,-1} = 0$):
\[
B \;=\; \sum_{j=0}^{n/2-1} 4^{j} \, d_j
\qquad\qquad
d_j \;=\; -\,2 \, B^{\,2j+1} + B^{\,2j} + B^{\,2j-1} \;\in\; \{-2, -1, 0, 1, 2\}
\]

\noindent No carry travels between digits: $d_j$ is a local function of its window.

\vspace{0.5em}

\noindent\textbf{Step 3 --- why it is exact}

\noindent Each $B^{\,2j+1}$ below the top is counted as $-2 \cdot 4^{j}$ in $d_j$ and
$+4^{j+1}$ in $d_{j+1}$: net $2^{\,2j+1}$, its weight in Step 1. The top bit is counted
once, as $-2 \cdot 4^{n/2-1} = -2^{\,n-1}$: its negative weight in Step 1. The sign is
built into the digit rule.

\vspace{0.5em}

\noindent\textbf{Step 4 --- the product}
\[
A \, B \;=\; \sum_{j=0}^{n/2-1} 4^{j} \, d_j A
\]

\noindent $d_j A$ takes one of five values:

\vspace{0.5em}

\begin{center}
\renewcommand{\arraystretch}{1.35}
\begin{tabular}{c c c | r | l}
$B^{\,2j+1}$ & $B^{\,2j}$ & $B^{\,2j-1}$ & $d_j$ & \multicolumn{1}{c}{$d_j A$} \\
\hline
$0$ & $0$ & $0$ & $0$  & $0$ \\
$0$ & $0$ & $1$ & $1$  & $A$ \\
$0$ & $1$ & $0$ & $1$  & $A$ \\
$0$ & $1$ & $1$ & $2$  & $2A$ \\
$1$ & $0$ & $0$ & $-2$ & $-2A$ \\
$1$ & $0$ & $1$ & $-1$ & $-A$ \\
$1$ & $1$ & $0$ & $-1$ & $-A$ \\
$1$ & $1$ & $1$ & $0$  & $0$ \\
\end{tabular}
\end{center}

\vspace{0.5em}

\noindent $n/2$ digits, $n/2$ partial products. Five digit values in radix four:
redundant, three control bits, full $\pm 2$ range, no carry.

\end{document}
